\documentclass[../main.tex]{subfiles}
\graphicspath{{\subfix{../assets/}}}

\begin{document}
\section{Brake Pad Coefficient of Friction}
\subsection{Calculating the coefficient}
Consider the definition for the coefficient of friction:
\[\mu \coloneq \dfrac{F_\text{max}}{N}\]
where $F_\text{max}$ is the maximum frictional force, and $N$ the normal force.
Hence, consider the following diagram.
\begin{figure}[h!]
	\begin{center}
		\includegraphics[width=0.55\textwidth]{incline.png}
	\end{center}
	\caption{Aiden's diagram of an incline plane}
\end{figure}
Hence, consider that as the system is static, $f = mg\sin\theta$, and that $\exists\theta:mg\sin\theta=F_\text{max}$.
Hence, the coefficient of friction is simply
\[\mu = \dfrac{\cancel{mg}\sin\theta}{\cancel{mg}\cos\theta} = \tan\theta \]
where $\theta$ is the angle of repose such that the object will tend to move. \\
Hence, we record the angle of repose across a number of trials:
\begin{table}[h!]
	\begin{tabularx}{0.5\textwidth} {
			| >{\centering\arraybackslash}X
			| >{\centering\arraybackslash}X |
		}
		\hline
		\textbf{Trial number} & \textbf{Measured $\theta$} \\
		\hline
		$1$                   & $28^\circ$                 \\
		\hline
		$2$                   & $31^\circ$                 \\
		\hline
		$3$                   & $26^\circ$                 \\
		\hline
		$4$                   & $28.5^\circ$               \\
		\hline
		$5$                   & $28^\circ$                 \\
		\hline
		$6$                   & $25^\circ$                 \\
		\hline\hline
		\textbf{Mean}         & $27.75^\circ$              \\
		\hline
	\end{tabularx}
\end{table} \\
And the result is a trivial calculation:
\[\mu = \tan(27.75^\circ) \approx 0.58. \tag{2DP}\]

\subsection{Determining a radius for the brake pad}
Consider that we want the average radius for the brake pad.
As it is nontrivial to directly take radius from the arc, consider the following diagram:
\begin{figure}[h!]
	\begin{center}
		\includegraphics[width=0.85\textwidth]{radiusdiagram.png}
	\end{center}
	\caption{Rex's geogebra diagram for calculating the radius}
\end{figure} \\
As we are given the arc centred at $A$ from $B$ to $C$, draw the line $BC$ and take $D$ as its midpoint. Draw the line $DE \perp BC$.
Consider, the radius is simply $AE$.
Define the following:
\begin{align*}
	r & \coloneq AE \\
	h & \coloneq DE \\
	c & \coloneq BC.
\end{align*}
Hence, by the Pythagorean theorem, we have:
\begin{align*}
	r^2          & = (r-h)^2 + \dfrac{c}{2}^2                  \\
	\cancel{r^2} & = \cancel{r^2} - 2rh + h^2 + \dfrac{c^2}{4} \\
	r            & = \dfrac{h}{2} + \dfrac{c^2}{8h}.
\end{align*}
And by a substitution of values, we have $r_\text{in}= 123$, $r_\text{out} \approx 170.8409$, and $r_\mu \approx 146.9205$, where $r_\mu$ denotes the arithmetic mean of the inner and outer radii.
\end{document}
